\documentclass{beamer}
%
% Volba vzhledu prezentace.
%
% Pro další možnosti, barevná schémata a použité fonty, navštivte:
% http://deic.uab.es/~iblanes/beamer_gallery/index_by_theme.html
%
\mode<presentation>
{
  \usetheme{default}      % nebo zkuste Darmstadt, Madrid, Warsaw,...
  \usecolortheme{default} % nebo zkuste albatross, beaver, crane, ...
  \usefonttheme{default}  % nebo zkuste serif, structurebold, ...
  \setbeamertemplate{navigation symbols}{}
  \setbeamertemplate{caption}[numbered]
} 

\usepackage[czech]{babel}
\usepackage[utf8]{inputenc}
\usepackage[T1]{fontenc}
\usepackage{lmodern}

%definice matematických prostředí
\newtheorem{veta}{Věta}
\newtheorem{lema}[veta]{Lemma}

\title[Ukázková prezentace]{Matematické modelovanie}
\author{Miroslava Vargová}
\institute{Techmická Univerzita v Košiciach}


\begin{document}

\begin{frame}
  \titlepage
\end{frame}


\begin{frame}{Obsah}
  \tableofcontents
\end{frame}

\section{Úvod}

\begin{frame}{Úvod}

\begin{itemize}
  \item Několik poznámek k tvorbě prezentací.
  \item Pro tvorbu výčtů použijte prostředí \texttt{itemize}.
\end{itemize}

\vskip 1cm

\begin{block}{Příklady}
Pro usnadnění práce při tvorbě prezentací uvádíme několik jednoduchých příkladů.
\end{block}

\end{frame}

\section{Příklady}

\subsection{Tabulky a obrázky}

\begin{frame}[fragile]{Tabulky a obrázky}

\begin{itemize}
\item Použijte prostředí \texttt{tabular} na sazbu základních tabulek -- viz. Tabulka~\ref{tab:widgets}.

\item Pro  vložení obrázku použijte příkaz \verb|\includegraphics|.
\end{itemize}

% Příkazy na vložení obrázku
\begin{figure}[htb]
\includegraphics[width=.4\textwidth]{mys.jpg}
\caption{\label{fig:mys}Bílá myška.}
\end{figure}

\begin{table}
\centering
\begin{tabular}{l|r}
Položka & Množství \\\hline
Jablka & 42 \\
Hrušky & 13
\end{tabular}
\caption{\label{tab:widgets}Ukázková tabulka.}
\end{table}

\end{frame}

\begin{frame}
\frametitle{Text a obrázek vedle sebe}
  \begin{columns}
    \begin{column}{.5\textwidth}
     \begin{block}{}
Lorem ipsum dolor sit amet, consectetur adipiscing elit. Vestibulum rutrum
urna vitae tellus consequat vestibulum. Nam ut purus diam, ac fermentum
metus. Sed quam augue, lobortis ac mattis non, consectetur vitae elit. Ut
tellus arcu, pellentesque convallis congue at, faucibus at velit. 
    \end{block}
    \end{column}%
    \begin{column}{.5\textwidth}
    \begin{block}{}
    \includegraphics[width=\textwidth]{mys.jpg}
    \end{block}
    \end{column}
  \end{columns}
\end{frame} 

\subsection{Matematika}\label{kapitola}

\begin{frame}[fragile]{Matematika}


\begin{veta}<1->
Moje první věta o tom, že 
$\sqrt{\vphantom{b} a} +\sqrt{b}$ 
vypadá lépe,
než $\sqrt{a} + \sqrt{b}$.
\end{veta}


\begin{veta}<2->[Abelova]
Věta s označením. 
\end{veta}

\begin{lema}<3->
$[0,1]$ je kompaktní podmnožinou $\mathbb R$.
\end{lema}
\end{frame}

\begin{frame}{Matematika -- pokračování}
Celý systém, skládající se z rovnic~\eqref{dcera1},
\eqref{dcera2} a~\eqref{dcera3}, se nachází 
v~Sekci~\ref{kapitola}.
\begin{alignat}{4}
  a_{11}x_1 &+ a_{12}x_2 &&+ a_{13}x_3 &&            &&=
y_1,\label{dcera1}\\
a_{21}x_1 &+ a_{22}x_2 &&            &&+ a_{24}x_4 &&=
y_2,\label{dcera2}\\
a_{31}x_1 &            &&+ a_{33}x_3 &&+ a_{34}x_4 &&=
y_3.\label{dcera3}
\end{alignat}

\begin{equation}
\cfrac{1}{2+\cfrac{1}{3 + \cdots} }\label{hvezdicka}
\end{equation}
Řetězové zlomky viz~\eqref{hvezdicka} nebo také~\cite[strana 30]{plch}.

\begin{equation}
\biggl\langle u \biggm| \sum_{i=1}^{n} F(e_i,v) e_i \biggr \rangle
= F \biggl ( \sum_{i=1}^n \langle  e_i | u \rangle e_i,v\biggr).
\end{equation}
\end{frame}

\section{Literatura}


\begin{frame}{Seznam použité literatury}

\begin{thebibliography}{9}
\bibitem{plch} PLCH, Roman; LOMTATIDZE, Lenka.
\emph{Sázíme v \LaTeX{}u diplomovou práci z~matematiky}. 1. vydání. Brno:
Masarykova univerzita, 2003. 122 s. ISBN 80-210-3228-6. 
\bibitem{rybicka} RYBIČKA, Jiří. \emph{\LaTeX{} pro začátečníky}. 3. vydání.
Brno: Konvoj, 2003. 238 s. ISBN 80-7302-049-1.
\end{thebibliography}
\end{frame}

\end{document}
